Krátká definice a účel: metakognice — student popíše vlastní myšlení, odůvodní postup a reflektuje učení; v hodnocení zařaďte pole pro sebereflexi.

Praktická doporučení:
\begin{itemize}
  \item Požadujte řešení + procesní artefakt (fotky/OneNote; OneNote převádí rukopisné rovnice) \cite{kb_ai_101_tipu_pdf_04e4a115_page_8_1}.
  \item Použijte automatizovanou předkontrolu (AI označí pravděpodobné chyby; při nízké jistotě eskalujte k učiteli).
  \item Vložte krátké metakognitivní otázky do zadání (např. „Popište 2 klíčová rozhodnutí a proč“).
\end{itemize}

Krátký workflow:
\begin{enumerate}
  \item Odevzdání řešení + artefakt.
  \item Předkontrola → při nejasnostech učitel.
  \item Student doplní krátkou reflexi; učitel provede spot‑check (5–10\%).
\end{enumerate}

Instrukce pro studenty: 
\begin{PromptBlock}

Nahrát: 1) výsledek, 2) krok-za-krokem (fotky/OneNote), 3) reflexe (max. 200 slov): co bylo klíčové? Kde nejistota?
\end{PromptBlock}
. Pozn.: Zvažte MS 365 Copilot vzhledem k GDPR \cite{kb_ai_101_tipu_pdf_04e4a115_page_3_1}.